% fixture: babel/switch-portuguese
% tier: T3
% engine: document
% covers: babel's \selectlanguage, otherlanguage, otherlanguage* and \foreignlanguage to a language other than the main one (english to portuguese and back): the patterns and hyphenmins of each paragraph, the \setlanguage whatsits in a paragraph, the captions and \today of the language, its shorthands ("-, "<) only while it is selected, the .aux writes
% issue: #338
\documentclass{article}
\usepackage[portuguese,english]{babel}
\pagestyle{headings}
\begin{document}
\section{Introduction}
Information theory studies the quantification, storage and communication of information in digital representations.
\begin{figure}[h]
\caption{An english figure}
\end{figure}

\selectlanguage{portuguese}
\section{Introdução}
A teoria da informação estuda a quantificação, o armazenamento e a comunicação da informação digital em representações. Palavras: guarda"-chuva e "<aspas">.
\begin{figure}[h]
\caption{Uma figura portuguesa}
\end{figure}
Com \foreignlanguage{english}{information representations} no meio, e mais transformações estabelecidas.

\selectlanguage{english}
Back in english, with \foreignlanguage{portuguese}{informação armazenamento guarda"-chuva} representations and "quoted" text,
\begin{otherlanguage*}{portuguese}transformações estabelecidas couve"-flor\end{otherlanguage*} words "here".
\begin{otherlanguage}{portuguese}
Um parágrafo em português com armazenamento e comunicações eletrônicas estabelecidas.
\begin{table}[h]
\caption{Uma tabela}
\end{table}
\end{otherlanguage}
Again english information representations.
\end{document}
